\chapter{Probas}
\label{chap:probas}

\lettrine{N}{este} capítulo explícase as probas que se realizaron durante o desenvolvemento da aplicación para comprobar o seu correcto funcionamento. Esta fase é moi importante xa que asegura que a implementación está libre de fallos e que ademais cumpre cos requisitos acordados con anterioridade.

\section{Probas unitarias}
As probas unitarias se realizan para comprobar o correcto funcionamento dunha unidade de código. Isto serve para asegurar que cada unidade funcione correctamente por separado. As probas realizáronse sobre os servizos da aplicación intentando cubrir os máximos casos posibles de forma que cada caso sexa independente do resto. Estas probas execútanse de maneira automática e verifícase que os resultados son os esperados. Tanto Ionic coma Django contan cun cliente que axuda a execución deste tipo de probas.

\section{Probas de sistema}
As probas de sistema teñen como obxectivo comprobar a integración de forma global, verifican que o funcionamento entre os distintos subsistemas que compoñen a aplicación é correcto. Estas probas realizáronse con dispositivos móbiles e estacións de traballo, utilizando en cada un os navegadores \textit{Google Chrome} e \textit{Firefox}. Non foi posible realizalas en dispositivos con sistemas operativos \textit{iOS} nin nos seus navegadores.

\section{Probas de aceptación}
Este tipo de probas realízanse unha vez finalizado todo o desenvolvemento, serven para determinar se a aplicación cumpre coas necesidades e os requisitos. Neste caso, as probas realizáronse polo autor e dous compañeiros que contaban con unha subscrición a plataforma de Spotify e comprobouse que cumpren os requisitos iniciais.
